\documentclass[12pt]{article}
\usepackage{amsmath,amssymb}

\begin{document}
\section*{{2020 AMC 12A Problems/Problem 6}}
Source: AoPS Wiki

\subsection*{Solution}
The two lines of symmetry must be horizontally and vertically through the middle. We can then fill the boxes in like so: [asy] import olympiad; unitsize(25); filldraw((1,3)--(1,4)--(2,4)--(2,3)--cycle, gray(0.7)); filldraw((0,3)--(0,4)--(1,4)--(1,3)--cycle, gray(0.9)); filldraw((3,3)--(3,4)--(4,4)--(4,3)--cycle, gray(0.9)); filldraw((4,3)--(4,4)--(5,4)--(5,3)--cycle, gray(0.9)); filldraw((2,2)--(2,3)--(3,3)--(3,2)--cycle, gray(0.9)); filldraw((2,1)--(2,2)--(3,2)--(3,1)--cycle, gray(0.7)); filldraw((3,0)--(4,0)--(4,1)--(3,1)--cycle, gray(0.9)); filldraw((1,0)--(2,0)--(2,1)--(1,1)--cycle, gray(0.9)); filldraw((4,0)--(5,0)--(5,1)--(4,1)--cycle, gray(0.7)); filldraw((0,0)--(1,0)--(1,1)--(0,1)--cycle, gray(0.9)); for (int i = 0; i < 5; ++i) { for (int j = 0; j < 6; ++j) { pair A = (j,i); } } for (int i = 0; i < 5; ++i) { for (int j = 0; j < 6; ++j) { if (j != 5) { draw((j,i)--(j+1,i)); } if (i != 4) { draw((j,i)--(j,i+1)); } } } [/asy] where the light gray boxes are the ones we have filled. Counting these, we get $\boxed{\textbf{(D) } 7}$ total boxes. ~ciceronii

\end{document}
